\begin{abstract}

\begin{revision}
Quantum constructions of Shor's order-finding on elliptic curves have established the feasibility of coherent group arithmetic and provided detailed resource estimates. Building on these results, this work presents a modular hybrid architecture that combines quantum order finding for the periodic function $f_P(k)=kP$, classical recovery and validation of the point order, and integration within an ECDSA-like signing workflow. The quantum stage follows the standard order-finding procedure and evaluates the periodic function through reversible elliptic-curve arithmetic based on established constructions. The contribution is twofold: an explicit integration of these components into a classical--quantum process, and a computational cost analysis that distinguishes controlled point addition from full scalar multiplication. The comparison with classical methods separates direct enumeration, generic group algorithms, and approaches based on point counting and factorization of the group order, making the assumptions and scope of the comparison explicit. Numerical validation combines a worked signing example with a finite reliability characterization of the implemented recovery policy, separating exact-order recovery, missing candidates, nonminimal multiples, and signing admission.

\end{revision}

\textbf{Keywords:} Quantum computing, elliptic curves, reversible arithmetic, Shor's algorithm, ECDSA, hybrid algorithms, quantum order finding.

\end{abstract}